\scanpage{78}
\mspage{71}
% archive sheet number: 77
est comp. avec les relations d'équivalences définies par
\[
\operatorname{Hom}(F',\unsure{G'|C'})\unsure{-|\longrightarrow}\operatorname{Hom}_{W(I')}(F',G'),
\]
resp
\[
\operatorname{Hom}(u_*F',u_*G')\longrightarrow\operatorname{Hom}_{W(I)}(u_*F',u_*G').
\]
Et il en est de $=$ de la flèche $\mu:f\mapsto f\circ i'$, trivialement. Donc la flèche d'adjonction $\operatorname{Hom}(u^*F,G')\longrightarrow\operatorname{Hom}(F,u_*G')$ est compatible avec les relations $R,R'$, ce qui achève de prouver que $(\unsure{\alpha\alpha|\alpha})$ page 68 induit une bijection $(\alpha)$ p. 68.

\deleted{\unsure{les}} Ceci montre donc $1^\circ)$ que pour un objet fixé $G'$ dans $\mathrm{HO}_W(I')$ (qu'on peut tjrs supposer, à isom près, provenir de $M(I')_f$) le foncteur
\[
F\longmapsto\operatorname{Hom}_{\mathrm{HO}_W(I')}(u_W^*F,G')
\]
est représentable --- i.e. le foncteur $u_W^*$ admet un adjoint à D., et $2^\circ)$ que cet adjoint à gauche se calcule comme $u_*(G')$
\margin{\unsure{Quel allonge\uline{ment}}}
($G'$ est pris fibrant) --- \deleted{comme} vu que l'hom $(\alpha\ \text{bis})$ page 69 est \uline{fonctoriel en $F$}. (A posteriori il s'ensuit qu'il est aussi fonctoriel en $G'$\ldots)
